\documentclass[11pt]{article}
\usepackage{../homework}

\profname{Dr. David Singer}
\classcode{MATH 465 Differential Geometry}
\assignnum{Assignment 07}
\duedate{02 November 2026}

\begin{document}
All numbered problems from Do Carmo, \textit{Differential Geometry of Curves and Surfaces} (2e).
\begin{problem}{1}
  The paraboloid of revolution \( M \) can be parametrized by \( \mathbf{x}(r,\theta) = (r\cos\theta, r\sin\theta, r^2) \).
  Compute the first fundamental form of \( M \) and the Gaussian curvature \( K(r,\theta) \) of \( M \).
\end{problem}
\begin{answer}
  First of all,
  \begin{equation}
  \label{eq:1}
  \begin{aligned}
    \mathbf{x}_r &= (\cos\theta, \sin\theta, 2r) \\
    \mathbf{x}_{\theta} &= (-r\sin\theta, r\cos\theta, 0).
  \end{aligned}
\end{equation}
Thus, the first fundamental form is
\begin{equation}
\label{eq:2}
ds^2 = (1+4r^2)dr^2 + r^2d\theta^2.
\end{equation}

Since the coordinates are orthgonal, we can use the Gauss formula,
\begin{equation}
\label{eq:3}
\begin{aligned}
  K &= - \frac{1}{2 \sqrt{EG}} \left[ \left( \frac{E_v}{\sqrt{EG}} \right)_v + \left( \frac{G_u}{\sqrt{EG}} \right)_u \right] \\
    &= - \frac{1}{2\sqrt{r^2 + 4r^4}} \left( \frac{2r}{\sqrt{r^2 + 4r^{4}}} \right)_{r} \\
    &= - \frac{1}{2r\sqrt{1+4r^2}} \left( \frac{2}{\sqrt{1+4r^2}} \right)_r \\
    &= - \frac{1}{2r\sqrt{1+4r^2}} \left( -\frac{8r}{(1+4r^2)^{3/2}} \right) \\
    &= \frac{4}{(1+4r^2)^2}.
\end{aligned}
\end{equation}
\end{answer}

\begin{problem}{2}
  Compute the total curvature \( \int\int_M K \dd{A} \) of \( M \).
  Note that \( \dd{A} = \sqrt{EG - F^2}\dd{r}\dd{\theta} \).
\end{problem}
\begin{answer}
  The area form is
  \begin{equation}
  \label{eq:4}
  \dd{A} = \sqrt{EG - F^2}\dd{r}\dd{\theta} = r \sqrt{1 + 4r^2}\dd{r}\dd{\theta}.
\end{equation}
Thus, the total curvature is
\begin{equation}
\label{eq:5}
\begin{aligned}
  \int\int_M K \dd{A} &= \int_0^{2\pi}\int_0^{\infty} \frac{4}{(1+4r^2)^2} r \sqrt{1 + 4r^2}\dd{r}\dd{\theta} \\
  &= 8\pi \int_0^{\infty} \frac{r}{(1 + 4r^2)^{3/2}}\dd{r}.
\end{aligned}
\end{equation}
Substitute \( y = 1 + 4r^2 \) so that \( \dd{y} = 8r\dd{x} \).
Then,
\begin{equation}
\label{eq:6}
\begin{aligned}
\int\int_M K \dd{A} &= \pi \int_1^{\infty} \frac{\dd{y}}{y^{3/2}} = -\pi\frac{2}{y^{1/2}}\eval_1^{\infty} = 2\pi.
\end{aligned}
\end{equation}
\end{answer}

\begin{problem}{4.4.5}
  Consider the torus of revolution generated by rotating the circle
  \begin{equation}
  \label{eq:7}
  (x-a)^2 + z^2 = r^2, \quad y = 0,
\end{equation}
about the \( z \) axis \( (a > r > 0) \).
The parallels generated by the points \( (a+r,0), (a-r,0), (a,r) \) are called the maximum parallel, minimum parallel, and upper parallel, respectively.
Check which of these parallels is
\begin{enumerate}[label=\alph*.]
  \item A geodesic.
  \item An asymptotic curve.
  \item A line of curvature.
\end{enumerate}
\begin{proof}
  We can parametrize a circle of radius \( r \) centered at \( a \) in the plane as \( x = r\cos\theta + a, z = r\sin\theta \), so that when we rotate about the \( z \) axis by an angle \( \varphi \) we get
  \begin{equation}
  \label{eq:9}
  \mathbf{x}(\theta,\varphi) = ((r\cos\theta + a)\cos\varphi, (r\cos\theta + a)\sin\varphi, r\sin\theta).
\end{equation}
The coordinate vectors are
\begin{equation}
\label{eq:10}
\begin{aligned}
  \partial_{\theta} &= (-r\sin\theta\cos\varphi, -r\sin\theta\sin\varphi, r\cos\theta) \\
  \partial_{\varphi} &= (-(r\cos\theta + a)\sin\varphi, (r\cos\theta + a)\cos\varphi, 0),
\end{aligned}
\end{equation}
giving the first fundamental form
\begin{equation}
\label{eq:11}
ds^2 = r^2d\theta^2 + (r\cos\theta + a)^2d\varphi^2.
\end{equation}

From this we can determine the connection coefficients
\begin{equation}
\label{eq:12}
\Gamma_{\mu\nu}^{\lambda} = \frac{1}{2}g^{\lambda\lambda} \left( \partial_{\mu}g_{\nu\lambda} + \partial_{\nu}g_{\mu\lambda} - \partial_{\lambda}g_{\mu\nu} \right).
\end{equation}
Hence,
\begin{equation}
\begin{aligned}
  \Gamma_{\theta\theta}^{\theta} &= 0, \quad \Gamma_{\theta\varphi}^{\theta} = 0, \quad \Gamma_{\varphi\varphi}^{\theta} = \frac{(r\cos\theta + a)\sin\theta}{r} \\
  \Gamma_{\varphi\varphi}^{\varphi} &=0, \quad \Gamma_{\theta\varphi}^{\varphi} = \frac{-r\sin\theta}{r\cos\theta + a}, \quad \Gamma_{\theta\theta}^{\varphi} = 0.
\end{aligned}
\end{equation}

Now, the points \( A = (a+r, 0), B = (a-r, 0), C = (a,r) \) can be seen to correspond to the planar polar coordinates \( \theta = 0, \pi, \pi/2 \), respectively.
Therefore, we can parametrize the generated curves as
\begin{equation}
\label{eq:13}
a^{\mu}(\lambda) = (0,\lambda), \quad b^{\mu}(\lambda) = (\pi, \lambda), \quad c^{\mu}(\lambda) = (\pi /2, \lambda),
\end{equation}
where \(\lambda \in [0, 2\pi)\).
For part a of the problem, we simply check whether these curves satisfy the geodesic equations
\begin{equation}
\label{eq:14}
\dv[2]{x^{\mu}}{\lambda} + \Gamma_{\alpha\beta}^{\mu} \dv{x^{\alpha}}{\lambda} \dv{x^{\beta}}{\lambda} = 0.
\end{equation}
The \(\mu = \varphi\) equation is easy, since \( \dv*[2]{\lambda}x^{\mu} = \dv*[2]{\lambda}(\lambda) = 0 \), and the only non-vanishing Christoffel symbol is \(\Gamma_{\theta\varphi}^{\varphi}\), but \( x^{\theta} = 0 \), so the whole equation vanishes.
For the \(\mu = \theta\) equation, again the second derivative term vanishes, but we do have the contribution from \(\Gamma_{\varphi\varphi}^{\theta}\).
However,
\begin{equation}
\label{eq:15}
\Gamma_{\varphi\varphi}^{\theta} \dv{x^{\varphi}}{\lambda}\dv{x^{\varphi}}{\lambda} = \frac{(r\cos\theta + a)\sin\theta}{r} = \frac{(r\cos(0) + a)\sin(0)}{r} = 0.
\end{equation}
Therefore, the maximum parallel is a geodesic.
Moreover, identical reasoning (since \(\pi \) is a constant and \( \sin(\pi) = 0 \)) shows that the minimum parallel is also a geodesic.
However, the upper parallel is not a geodesic, since the terms involving \(\Gamma_{\varphi\varphi}^{\theta}\) are proportional to \( \sin(\pi/2) = 1 \).

To check asymptotics, we will need to calculate the coefficients of the second fundamental form.
The derivatives of the coordinate vectors are
\begin{equation}
\label{eq:16}
\begin{aligned}
  \partial_{\theta\theta} &= (-r\cos\theta\cos\varphi, -r\cos\theta\sin\varphi, -r\sin\theta) \\
  \partial_{\theta \varphi} &= (r\sin\theta\sin\varphi, -r\sin\theta\cos\varphi, 0) \\
  \partial_{\varphi\varphi} &= (-(r\cos\theta + a)\cos\varphi, -(r\cos\theta + a)\sin\varphi, 0),
\end{aligned}
\end{equation}
a surface normal is
\begin{equation}
\label{eq:17}
\begin{aligned}
  n &= \partial_{\theta} \wedge \partial_{\varphi} = \left( -r\cos\theta(r\cos\theta + a)\cos\varphi, -r\cos\theta(r\cos\theta+a)\sin\varphi, -r\sin\theta(r\cos\theta + a) \right)
\end{aligned}
\end{equation}
which has norm
\begin{equation}
\label{eq:18}
|n| = r(r\cos\theta + a),
\end{equation}
so that a unit normal is
\begin{equation}
\label{eq:19}
\nu = (\cos\theta\cos\varphi, \cos\theta\sin\varphi, \sin\theta).
\end{equation}
Finally,
\begin{equation}
\label{eq:20}
\begin{aligned}
  e &= \langle \partial_{\theta\theta}, \nu \rangle = -r, \\
  f &= \langle \partial_{\theta\varphi}, \nu \rangle = 0 \\
  g &= \langle \partial_{\varphi\varphi}, \nu \rangle = -\cos\theta(r\cos\theta + a).
\end{aligned}
\end{equation}
We now check if any of the parallels satisfy
\begin{equation}
\label{eq:21}
\mathrm{I \!I}(\dv{x}{\lambda}) = 0 = e \dot{x^{\theta}}^2 + 2f(\dot{x^{\theta}}\dot{x^{\varphi}}) + g \dot{x^{\varphi}}^2.
\end{equation}
In fact, all of the parallels have tangent vectors \( \dot{a}^{\mu} = \dot{b}^{\mu} = \dot{c}^{\mu} = (0,1) \).
Therefore, we calculate
\begin{equation}
\label{eq:22}
\begin{aligned}
  \mathrm{I \!I}(\dot{a}^{\mu}) &= \cos\theta(r\cos\theta + a) = -(r + a) \neq 0 \\
  \mathrm{I \!I}(\dot{b}^{\mu}) &= -(-r + a) \neq 0 \\
  \mathrm{I \!I}(\dot{c}^{\mu}) &= 0.
\end{aligned}
\end{equation}
Therefore only the upper parallel is an asymptotic curve.

Finally, since both the first and second fundamental forms are diagonal, we can find the Weingarten map easily as
\begin{equation}
\label{eq:23}
L_i^j = g^{jk}L_{ik} = \begin{pmatrix}
- \frac{r}{r^2}  &  \\
                 & - \frac{\cos\theta(r\cos\theta + a)}{(r\cos\theta + a)^2} \end{pmatrix} = \begin{pmatrix}
- \frac{1}{r}  &  \\
               & - \frac{\cos\theta}{(r\cos\theta + a)}
 \end{pmatrix}.
\end{equation}
This gives the principal curvatures
\begin{equation}
\label{eq:24}
\kappa_1 = - \frac{1}{r}, \quad \kappa_2 = - \frac{\cos\theta}{(r\cos\theta + a)},
\end{equation}
which clearly correspond to principal directions \( e_1 = (1,0), e_2 = (0,1) \), or in other words, the coordinate curves.
Each of the parallel curves have tangent vector \( (0,1) \), meaning all three are eigenvectors of the Weingarten map, so all three are lines of curvature.
\end{proof}
\end{problem}

\begin{problem}{4.4.6}
  Compute the geodesic curvature of the upper parallel of the torus of Exercise 4.4.5.
\end{problem}
\begin{answer}
  We know that the geodesic and normal curvatures satisfy
  \begin{equation}
  \label{eq:29}
  k^2 = k_g^2 + k_n^2,
\end{equation}
and that we can get the normal curvature as \( \mathrm{I \!I}(T) \) where \( T \) is the unit tangent vector.
\end{answer}

\begin{problem}{4.5.3}
  Let \( S \subset \mathbb{R}^3 \) be a regular compact surface with \( K > 0 \). Let \( \mathbf{\Gamma} \subset S \) be a simple closed geodesic in \( S \), and let \( A \) and \( B \) be the regions of \( S \) which have \( \mathbf{\Gamma} \) as a common boundary.
  Let \( N: S \rightarrow S^2 \) be the Gauss map of \( S \).
  Prove that \( N(A) \) and \( N(B) \) have the same area.
\end{problem}
\begin{proof}

\end{proof}
\end{document}
